\documentclass[11pt]{article}

\usepackage[margin=0.75in]{geometry}
\usepackage{booktabs}
\usepackage{amsmath}
\usepackage{multirow}
\usepackage{hyperref}

\title{Dual-Domain MMD Adaptation: Held-Out Test Set and Large-Target Follow-Up}
\author{}
\date{\today}

\begin{document}
\maketitle

\section{What this document adds}

Two things not present in \texttt{experiment\_summary\_v6.tex} or
\texttt{freeze\_epoch\_ablation.tex}: (1) a genuinely held-out real test set, and
(2) a first look at \texttt{real\_large} (rather than \texttt{real\_small}) as the
adaptation target from a \texttt{syn\_xlarge} source, at both YOLOv10n and YOLOv10m
scale.

\subsection{The held-out test set}

Every real-domain comparison in this study up to now, baselines included, has been
scored on the same 548-image real validation set -- which is also the set used for
checkpoint selection (\texttt{best.pt}) and early stopping during training. That is a
mild form of leakage: the val set is never trained on directly, but it does influence
which epoch's weights get kept.

\texttt{real\_test\_heldout} closes that gap. It is the official VisDrone2019-DET
\emph{test-dev} split (1610 images), which is a separate download from the
train/val archive this project's \texttt{real\_small}/\texttt{real\_large} both draw
from (\texttt{scripts/build\_domain\_datasets.py}'s \texttt{DEFAULT\_REAL\_ROOT}).
Filename-level comparison confirmed zero overlap with the train+val pool used by any
run in this study. Labels were remapped from VisDrone's raw 12-category scheme to
this project's person/vehicle scheme, bicycles and ignored-regions dropped, matching
the \texttt{yolo\_splits\_no\_bikes} convention already used for \texttt{real\_small}/
\texttt{real\_large}. All 1610 images retained at least one box after the remap.
Built non-destructively: images symlinked, nothing existing modified.

\section{real\_small target, re-checked on the held-out set}

\begin{table}[h!]
    \centering
    \begin{tabular}{lrrrr}
        \toprule
        Model & person mAP50 & person mAP50-95 & vehicle mAP50 & vehicle mAP50-95 \\
        \midrule
        baseline (real\_small, solo) & 0.433 & 0.178 & 0.713 & 0.463 \\
        MMD (regularized, syn\_xlarge pretrain) & 0.403 & 0.163 & 0.699 & 0.451 \\
        \bottomrule
    \end{tabular}
    \caption{\texttt{real\_small\_1920\_testing-3} (solo, COCO-direct) vs.\
    \texttt{syn\_xlarge\_yolo\_n\_source\_real\_small\_target\_reg\_fr5}, both
    evaluated on \texttt{real\_test\_heldout} (1610 images, never used in any
    train/val split in this study).}
\end{table}

\noindent Solo wins on every column, slightly more decisively than on the shared
\texttt{real\_large} val set used elsewhere in this study -- ruling out shared-val
checkpoint-selection leakage as the explanation for that earlier result.

\section{real\_large as the adaptation target (new)}

Source = \texttt{syn\_xlarge} (YOLOv10n pretrain, 150 epochs, imgsz 1920), target =
\texttt{real\_large} (6470 train images, vs.\ \texttt{real\_small}'s $\sim$1600) --
the first time this study has used \texttt{real\_large} rather than
\texttt{real\_small} as the MMD target. Two MMD variants are compared against a solo
(COCO-direct, no synthetic detour) baseline trained on the same target set.

\begin{table}[h!]
    \centering
    \begin{tabular}{lrrrr}
        \toprule
        Model & \multicolumn{2}{c}{real\_large val} & \multicolumn{2}{c}{real\_test\_heldout} \\
        & mAP50 & mAP50-95 & mAP50 & mAP50-95 \\
        \midrule
        real\_large\_1920 (solo, COCO-direct) & 0.805 & 0.505 & 0.617 & 0.349 \\
        syn\_xlarge\_source\_real\_large\_target\_reg\_fr5 (freeze=6) & 0.790 & 0.496 & 0.598 & 0.339 \\
        \quad{}\_stronger\_weight-5 (120ep, freeze=25) & \textbf{0.802} & \textbf{0.501} & \textbf{0.606} & \textbf{0.342} \\
        \bottomrule
    \end{tabular}
    \caption{\texttt{real\_large} as MMD target, source \texttt{syn\_xlarge}
    (YOLOv10n). Solo wins on both eval sets against both MMD variants, but the
    stronger-weight variant narrows the gap to the tightest margin seen anywhere in
    this study (0.003--0.007 mAP50-95).}
\end{table}

\noindent The matched no-MMD control for this regime (plain fine-tune of the
\texttt{syn\_xlarge} checkpoint on \texttt{real\_large}, no MMD term) has not yet been
run. Without it, the gap above conflates two variables -- the synthetic-pretrain
detour itself, and the MMD term -- and cannot yet be attributed to MMD specifically.
This is the natural next run to close out this section.

\subsection{Model scale: YOLOv10m}

\begin{table}[h!]
    \centering
    \begin{tabular}{lrrrr}
        \toprule
        Eval set & person mAP50 & person mAP50-95 & vehicle mAP50 & vehicle mAP50-95 \\
        \midrule
        real\_large & 0.679 & 0.324 & 0.894 & 0.647 \\
        real\_test\_heldout & 0.446 & 0.186 & 0.726 & 0.479 \\
        \bottomrule
    \end{tabular}
    \caption{\texttt{syn\_xlarge\_source\_real\_small\_target} at YOLOv10m scale
    (source = \texttt{syn\_xlarge\_yolo\_m-5}, target = \texttt{real\_small}).
    Beats every YOLOv10n variant in this study on both eval sets, as expected for
    the larger backbone. No YOLOv10m solo baseline exists yet to isolate MMD's
    contribution at this scale.}
\end{table}

\section{The other direction still holds: real$\to$syn}

\begin{table}[h!]
    \centering
    \begin{tabular}{lrrrr}
        \toprule
        Model & person mAP50 & person mAP50-95 & vehicle mAP50 & vehicle mAP50-95 \\
        \midrule
        syn\_small\_1920\_testing-2 (solo) & 0.667 & 0.332 & 0.684 & 0.465 \\
        real\_large\_source\_syn\_small\_target\_reg-4 (MMD) & \textbf{0.678} & \textbf{0.344} & \textbf{0.691} & \textbf{0.477} \\
        \bottomrule
    \end{tabular}
    \caption{Source = \texttt{real\_large\_1920}, target = \texttt{syn\_small},
    evaluated on \texttt{syn\_large} val (640 images). This is the only comparison
    in this document where MMD beats solo on every column -- consistent with
    \texttt{experiment\_summary\_v1}/README already flagging real$\to$syn as the
    direction where adaptation genuinely helps, unlike syn$\to$real.}
\end{table}

\section{Diagnostic: bandwidth-freeze does not prevent late-training drift}

Investigated why \texttt{train/mmd\_distance} climbs back up after the bandwidth
freeze, across two runs with different freeze epochs:

\begin{itemize}
    \item \textbf{freeze=6}: distance drops fast (177{,}446 $\to$ 157) through epoch
    6, then climbs straight back to $\sim$25{,}000 by epoch 11 -- superficially
    consistent with freezing mid-collapse, before the EMA had stabilized.
    \item \textbf{freeze=25}: distance was already low and fluctuating
    (18--1268) for several epochs \emph{before} freezing, i.e.\ genuinely stabilized
    -- yet still climbed steadily afterward (256 at epoch 27 $\to$ 7870 by epoch 45,
    roughly 300--400/epoch). This rules out ``froze too early'' as the explanation.
\end{itemize}

\noindent The behavior looks structural rather than a tunable timing issue: once
frozen, the bandwidth stops tracking current feature geometry, and the MMD weight is
simultaneously decaying by design (to avoid the collapse feedback loop documented in
\texttt{experiment\_summary\_v2.tex}). Once detection loss dominates enough, nothing
anchors the backbone's features to where they were at freeze time, so raw distance
from that now-stale reference drifts upward for the remainder of training regardless
of which epoch the freeze happened at. Detection mAP kept improving normally in both
affected runs, so this reads as an expected loss of signal in a late-training
diagnostic rather than a training pathology worth chasing further via freeze-epoch
tuning.

\section{Open items}

\begin{itemize}
    \item No-MMD control for \texttt{real\_large} target (source \texttt{syn\_xlarge},
    plain fine-tune, no MMD term) -- needed to attribute the real\_large-target gap
    to MMD specifically rather than the synthetic-pretrain detour alone.
    \item No YOLOv10m solo baseline on \texttt{real\_small} or \texttt{real\_large} --
    needed to check whether the solo-beats-MMD pattern holds at this model scale too.
    \item \texttt{real\_test\_heldout} has only been used with real-domain,
    \texttt{real\_small}/\texttt{real\_large}-target checkpoints so far; no equivalent
    held-out set exists yet for the synthetic domain.
\end{itemize}

\end{document}
